\section{La fusión de la simulación física y la IA}

\subsection{La brecha entre los tres mundos}

En la entrevista hay un pasaje excelente sobre los «tres mundos»: (1) el mundo real; (2) el mundo que podemos observar a simple vista; (3) el mundo que modelamos con modelos físicos.

\begin{importantbox}{El Gap entre los tres mundos}
\textbf{El primer Gap} (mundo real vs. mundo observable): cada persona solo puede observar una parte ínfima del mundo, así que el conocimiento de cada persona tiene limitaciones enormes. La esencia de este gap es \textbf{lo limitado de la capacidad de observación}.

\textbf{El segundo Gap} (mundo observable vs. mundo simulado): «hay una distancia enorme». Pero la razón es distinta: este gap existe porque \textbf{no tenemos suficiente compute ni suficiente data} para modelar el mundo real. En algunos escenarios clásicos (como la correlación entre CFD y las pruebas de túnel de viento) ya es muy buena, pero «antes de que un avión despegue de todos modos hay que probarlo en el túnel de viento, no se puede decir que con el CFD ya basta».
\end{importantbox}

Su reflexión al respecto está cargada de sentido filosófico: «Es bastante desalentador. Tal vez nosotros, como especie, vivamos dentro del simulador de alguien más, y aun así fantaseamos con construir nuestro propio simulador capaz de simularnos a nosotros mismos». Pero al mismo tiempo considera que un modelo como GPT-4, capaz de mimic el funcionamiento del cerebro humano, ya «es un milagro».

\subsection{Video Model vs. Physical Simulation}

«¿Por qué seguimos usando Taichi para hacer rendering, para hacer physical simulation? ¿Por qué no usar Sora?» — esta pregunta apunta directo al valor de la existencia de la disciplina de Graphics en la era de la IA.

Hu Yuanming (胡渊鸣) respondió con mucho pragmatismo: «Hace cinco años, Frozen se escribía en 99 líneas; ahora, al prompt que se le da a Sora le basta con 20 palabras. En algunos escenarios de verdad se le acabó el oficio a la gente de Graphics. Pero muchas de las cosas que hace Graphics fueron las que hicieron posible que naciera Sora.»

La dirección final de esta evolución es «predominantemente learning», pero el simulator sobrevivirá de alguna forma: como inductive bias, como fuente de synthetic data, o como parte de un learnable module. Neural Graphics (DLSS, NeRF, etc.) ya está reemplazando gradualmente al ray tracing tradicional, y el posicionamiento de Taichi 2.0 es precisamente «supporting Graphics and AI».

\subsection{Sim-to-Real Gap: el reto fundamental de la simulación}

Hu Yuanming señala que todos los simulator existentes (Isaac, MuJoCo, Drake, PyBullet, Genesis) enfrentan un problema en común: el \textbf{sim-to-real gap}.

\begin{warningbox}{Las leyes físicas son exactas, pero las condiciones de frontera no}
La segunda ley de Newton es exacta, y la relatividad también lo es. Pero muchas boundary condition no lo son: por ejemplo, el modelo de fricción; el mecanismo de fricción de una superficie lisa y el de una superficie rugosa son completamente distintos, y la fricción de un vidrio sobre una superficie mojada es distinta otra vez. Todos estos modelados son inexactos.
\end{warningbox}

Él considera que los simulator del futuro necesitan volverse más data-driven y combinarse con redes neuronales. Esto no significa usar un Transformer para volver a aprender la segunda ley de Newton (no tiene sentido gastar miles de millones de FLOPS en algo que originalmente se puede hacer con unos cuantos FLOPS), sino dejar que los métodos data-driven compensen las partes que \textbf{los modelos aproximados no pueden describir con precisión}, como la stress-strain curve de los materiales.

\subsection{Inductive Bias: el valor del prior físico en la IA}

Hu Yuanming subrayó en particular la importancia del inductive bias (sesgo inductivo) en los sistemas de IA:

\begin{importantbox}{La eficiencia del 3D Inductive Bias}
Tomando la video generation como ejemplo: en la graficación tradicional, la camera projection solo necesita unos 50 FLOPS, pero si se le pide a la IA que aprenda ese mapeo, tal vez haya que agregar parámetros del orden de billion. Por eso, la camera projection debería quedar como hardcode, dejándole al modelo la capacidad para las partes que de verdad necesitan aprenderse. De forma similar, en un video model, cuando la camera gira 360 grados aparece un mundo completamente distinto; si existe una representación 3D básica, al menos se puede resolver el problema de physical persistency.
\end{importantbox}

Resumió esto en un principio de equilibrio muy incisivo:

\begin{quote}
«El simulator necesita tener más AI adentro, y la AI necesita tener más 3D inductive bias adentro. Hay que encontrar el punto de equilibrio exacto: no hay que usurpar funciones que no corresponden, no hay que subestimar la capacidad de aprendizaje del modelo, pero tampoco hay que dejarle todo al learning.»
\end{quote}

\subsection{Las limitaciones de Synthetic Data}

Sobre la resistencia de Sergey Levine hacia la simulation en su blog «The Spork of AGI» —donde sostiene que la simulation en realidad es una limitación para los robotics foundation model y que se debería usar real world data—, Hu Yuanming dijo que «tiene sentido»:

\begin{warningbox}{Synthetic Data no puede ser la fuente principal de datos}
«Si toda la data es synthetic data, o si la mayor parte de la data es synthetic data, definitivamente no va a work. Porque la data que se puede synthesize siempre tiene que follow ciertas rules, y esas rules necesariamente las pensó un ser humano. De cien millones de situaciones en el mundo real, las rules tal vez solo alcancen a cover diez mil.» — Pero synthetic data sí puede usarse para enhance un modelo ya existente, en lugar de build from scratch.
\end{warningbox}

\subsection{Los límites de la simulación y los universos anidados}

Al hablar sobre «los límites de la simulación física», Hu Yuanming dio una respuesta que invita a la reflexión. Si nuestro propio mundo ya es algo simulado, entonces \textbf{el simulador dentro del simulador no puede ser más potente que el simulador matriz}.

\begin{knowledgebox}{El argumento de complejidad computacional de la simulación anidada}
«Tal vez tú, dentro de nuestro mundo, construyas un simulador que simule el mundo entero y que se ejecute más rápido que el tiempo real; eso probablemente sea unlikely. A menos que solo hagas una simulación small scale. Esto también significa que \textbf{predecir el mercado de valores probablemente sea muy difícil}, porque el mercado de valores involucra a todas las personas del mundo entero participando en él: el sistema que necesitarías simular sería tan grande como el mundo real mismo. Pero hacer una simulación de un túnel de viento sí es posible.»
\end{knowledgebox}

Esta perspectiva conecta con elegancia la reflexión filosófica de la infancia sobre «si el árbol de durazno crece cuando no lo estoy viendo» con la teoría de la complejidad computacional de la vida adulta.

\subsection{La caverna de Platón de los LLM}

Cuando le preguntaron si «los LLM como GPT viven en la caverna de la que hablaba Platón», la respuesta de Hu Yuanming fue inesperada:

\begin{quote}
«Sí. Porque su loss function es su training data, y todo lo que hace es minimize esa loss function. Pero, ¿acaso cada uno de nosotros no vive también dentro de la caverna de Platón? Es posible que cada uno de nosotros vea incluso menos que un LLM. Solo que el LLM a veces comete errores de sentido común, y por eso uno se ríe de él seguido. Pero, ¿de verdad somos mejores que el LLM? Yo creo que no necesariamente.»
\end{quote}

Él considera que la AGI puede surgir a partir de la text, video e image data ya existente, y que no necesariamente hace falta interiorizar en ella todo el conocimiento humano sobre el mundo físico. Lo esencial es darle a la AGI la capacidad de \textbf{tool use}: dejar que tenga su propia computadora y que pueda escribir código para analizar fenómenos físicos. «Puede descubrir por sí misma su propia ley de Newton.»

\subsection{Si se hiciera de nuevo un Simulator}

Si hoy se volviera a hacer un simulator, Hu Yuanming lo impulsaría \textbf{predominantemente con learning}, y no con la ley de Newton como base. Los «elementos subyacentes» que conservaría incluyen:
\begin{itemize}
\item Los objetos no desaparecen de la nada (physical persistency)
\item Dos objetos no pueden ocupar la misma posición (collision)
\item La perspective projection de la camera
\end{itemize}

Y convertiría todas las physical law (incluidas la deformación elastoplástica, la dinámica de fluidos, etc.) en algo data-driven. El sistema resultante no sería ni un simulator tradicional ni una neural network pura, sino \textbf{«un 3D inductive bias con una neural network incrustada»}: un prior estructurado que le facilita el aprendizaje a la neural network.

\subsection{¿La AGI necesita física?}

Hu Yuanming considera que la AGI «tal vez no necesite en particular la ley de Newton ni las ecuaciones de Lagrange», pero sí necesita \textbf{tool use}: la AGI necesita tener su propia computadora, capaz de escribir código finite element para analizar fenómenos físicos.

\begin{knowledgebox}{La limitación esencial del Next Token Prediction}
Sin importar cuánta o cuán poca información contenga un token, cada token de un Transformer consume una cantidad fija de compute ($2 \times$ parameter count / sparsity via MoE). Pero, por naturaleza, algunos token necesitan más compute que otros; ese compute adicional no puede ocurrir dentro del paradigma actual del LLM, sino que requiere que el LLM opere herramientas externas. Por eso la AGI no necesita que se le interiorice todo el conocimiento humano sobre el mundo físico: puede «descubrir por sí misma su propia ley de Newton».
\end{knowledgebox}

Esta idea tiene una consecuencia importante: la forma actual de scaling up un Transformer (aumentar la cantidad de parámetros, para que cada token reciba más compute) es, en esencia, resolver por fuerza bruta un problema que existe a nivel de la arquitectura. Un enfoque más inteligente tal vez sea —como hace Taichi— procesar de manera distinta cómputos de granularidades distintas.

\subsection{Reflexiones sobre Robotics}

Aunque la formación técnica de Hu Yuanming está muy relacionada con Robotics (simulación física, programación diferenciable, generación 3D), él admite que su «nivel de interés en Robotics no es tan alto».

\begin{quote}
«Creo que Robotics es algo con un impact enorme. Tal vez la única razón por la que no me he puesto a hacer Robotics es que, hoy por hoy, mi nivel de interés en esto no es tan alto. Creo que solo puedo hacer bien las cosas que de verdad me gustan.»
\end{quote}

Pero también reconoce que existe un punto de intersección potencial entre Robotics y Meshy: muchos socios ya están usando Meshy para generar activos 3D como tazas, tazones y mesas para el entrenamiento de Robotics. Si en el futuro Meshy soporta articulated objects (objetos con articulaciones), esta intersección se volverá aún más notable.

Su juicio sobre el papel de la simulation en Robotics también es muy perspicaz: hacer que un robot manipule las verduras y los platos sobre una mesa \textbf{es más difícil que hacer volar un avión}: «porque allá arriba no hay más que aire, pero cuando necesitas manipulate algo sobre la mesa con más detalle, la inconsistencia entre las distintas condiciones de frontera hace que sea muy difícil llevar la simulation a la realidad.»

\subsection{Resumen de la sección}

La simulación física atraviesa un momento clave de transición de rule-based a data-driven. El cuello de botella de la simulación tradicional no está en las leyes físicas en sí, sino en el modelado preciso de las condiciones de frontera. La dirección futura es conservar el prior 3D subyacente (como inductive bias) y dejar que el learning se encargue de las partes complejas y difíciles de modelar con precisión. Taichi 2.0 está evolucionando en esa dirección.
